Moving JHQ to a GPU raises three technical challenges.

First, compressed codes reduce storage but can still be expensive to scan.
Each eight-bit primary code addresses a 256-entry distance lookup table
(LUT); these tables must be built and may be too large for fast on-chip memory.
Because JHQ combines per-coordinate quantisation levels, this table can
be split exactly in real arithmetic into two 16-entry tables, one for
each four-bit half of the byte. The challenge is whether the smaller
tables repay the extra lookups and additions on the GPU.

Second, small batches can contain too few queries to keep the GPU busy.
Splitting each query's candidates across thread blocks creates parallel
work, but equal local quotas can discard global winners if the best
candidates fall in one block. One solution is to let each group retain
its own full shortlist before merging, at the cost of memory and merging
work. The challenge
is to retain all potential winners while keeping this overhead below
the parallelism gain. This thesis studies shared cutoffs and batch-dependent
scheduling as ways to reduce that overhead.

Third, the refinement work needed for a target recall varies by workload.
Too few candidates can discard true neighbours; too many spend time
reading and scoring residual codes without improving the answers.
The aim is to choose a reusable budget without ground-truth labels.
A larger-budget search provides a reference, but can itself miss true
neighbours. Agreement with it therefore needs separate checks of recall
on test queries, alongside the cost of calibration and subsequent search.
